因此，问题二更像是在搭建一个“珍稀物种响应模块”。题面附件中的事实需要逐一映射到参数：2022 年公报对应种群初值和监测基准，附件 1 中的航道阻隔、鱼道连通与栖息地破碎化对应 \(\delta_B\)，附件 2 中的放流、补充和食源恢复对应 \(R_S\)，而 \(e_D,m_D,H,u\) 则分别反映繁殖效率、自然死亡、捕捞/扰动和污染压力。只有把这些事实放进统一框架，才能说明禁渔收益是如何在食物网和生境连通性上被放大或削弱的，也才能解释为什么同处一个背景下，江豚和长江鲟会表现出不同的恢复曲线\cite{report2022}。

题面中的事实映射还意味着，问题二必须把“自然恢复”和“工程补偿”分开解释。对于江豚，更多是食源改善后的数量响应；对于长江鲟，除了食源，还要看洄游通道是否连续、放流是否提供额外补充、污染是否抵消恢复收益。因此，本问的重点不是比较哪个物种的绝对值更大，而是说明政策收益在不同约束条件下被保留了多少、损失了多少。把这层关系写清楚，才能让问题二从“珍稀物种是否增加”上升为“珍稀物种响应机制如何分化”。

表格的作用也在于让读者一眼看出每个参数到底来自哪条事实链，而不是把变量当成纯数学符号使用。这样写，问题二的证据链就会更完整：上游是禁渔和食源恢复，中间是通道、放流与污染修正，下游才是珍稀物种的数量变化。

